\subsection{A complete 198-state implementation of the new gate}
\label{sec:distributed-machine}

The shortest fully checked implementation of the distributed gate produced
in this investigation has \textbf{198 working states}. This is a concrete
compilation result, but it does not improve the existing 196-state bound.
The directory \file{distributed198/} contains the full machine, its exact
source tree and compiler, the selected macro word, and a portable verifier.

The dense table \file{distributed198/candidate.tm} has working states
$0,\ldots,197$, initial state 0, and a separate halt label 198. There are
396 transition rows. Its SHA-256 is
\begin{quote}\small
\nolinkurl{316c1fd0ade319e3896037f969c18065ecb690339816a5e1defb797ca6751b0c}.
\end{quote}
The native table uses initial state 74, interpreter root 84, and halt 254;
\file{labels.json} gives the complete conjugacy. All 198 working states
are reachable in the rooted transition graph. No claim that all are visited
on the actual blank computation is needed for the count.

Its physical budget is
\[
 134\text{ unchanged controls}+67\text{ AX readers}
        -3\text{ contextual folds}=198.
\]
The existing native reader 99 supplies the two-increment subtree. The three
folds are applied simultaneously and their complete transition edits are
recorded in \file{folds.json}. They retain the original seven-zero scanner.
The maximum source address length is 37.

The raw 201-state precursor and the folded 198-state table both pass all
4,908 complete symbolic instruction contracts for the revised 3,773-site
program. The folded check uses 1,585,320 symbolic execution steps in total,
at most 497 per contract. The arbitrary-gap and whole-register-file
interpretation is the same as for the 196-state rotations: all rows used
by the inherited register-block contracts are checked unchanged. These
counts are symbolic execution counts, not elapsed steps of the blank run.

The literal blank startup takes exactly 58,154 tape transitions. It reaches
dense root 75 (native 84), head position $-2$, source site 25, and the same
complete clean frame as PA196, including gap 48 and all 26 spare one-blocks.
This equality includes the entire finite tape support and all zeros outside
it. The finite startup is used only to establish the simulation's initial
condition.

\begin{lemma}[The selected emitter's nine stack entries]
Let the incoming accumulator be $t$, the incoming variable registers
be $d,h$, and the three post-prefix parameters be $D,H,A$.
After the selected AX word, the nine newly pushed entries, from top
to bottom, are
\[
 h5,\ h4,\ h3,\ h2,\ h6',\ h1,\
 \rho8(D,A,H,d),\ U(d,h),\ t,
\]
where $h6'$ is the antecedent-exchanged distribution head.
The enabled accumulator is the distributed gate $F_*(H,D,A)$;
the disabled accumulator is zero under the unchanged guard.
\end{lemma}
\begin{proof}
For a constructor expression $X$, its mathematical postfix emitter
takes an incoming accumulator $v$ and stack $w$ to accumulator $X$
and stack $v::w$. For a variable or constant-one expression this is
the copy/push definition. In the binary-constructor case, the two
recursive emissions push the incoming accumulator and then the
first child; the constructor consumes the latter, leaving exactly
the former. Tag increments do not change the stack. This proves
the claim by structural induction on $X$.

The unchanged 18-token prefix ends with accumulator
$\rho8(D,A,H,d)$ and stack prefix $[U(d,h),t]$. The six successive
head emissions have order $[1,6',2,3,4,5]$; the final gate emission
pushes the last head. Applying the preceding stack identity gives
exactly the displayed list. The unchanged guard suffix preserves
that stack and either keeps the gate or clears the accumulator.
\end{proof}

An additional independent checker,
\file{axiom_rewrite/audit_physical_binding.py}, imports neither the compiler
nor the symbolic runtime. It parses the entire delivered source tree,
including its root; proves that it is the pinned PA196 source with only
the AX subtree replaced; reconstructs the selected 199-token word; resolves
every primitive control edge by subtree intervals; and compares all 3,773
site records with the emitted certificate. It checks every macro occurrence
against the pinned macro's relative control graph, the 169 exterior sites
and their relocated edges, and the final closing jump. It also independently
replays the full blank startup. This audit passed without a discrepancy.

The logical and physical interfaces are now explicit: the machine implements
the revised emitter, and the revised numerical calculus has the PA-verifiable
proof translations of Theorem~\ref{thm:distributed-pa-translation}. The
ordinary source-level nonhalting transfer uses the inherited scheduler and
proof-search argument with the updated head-selection ranks. This package
does not export a new Lean build of those arguments or a complete internal
PA theorem about the new literal halting sentence. In particular, its
relative numerical-calculus theorem does not fill the separately documented
baseline consistency-bridge obligations.
